\section*{Limitations}
\label{sec:limitations}

\fablenote{Every limitation below is immediately neutralized with a ``but this is easy to fix'' clause, which reads as spin. A+ practice: state at least one limitation plainly and without rebuttal --- the honest one here is the 10-point average accuracy gap on general VQA and the fact that only 24--30\% of queries are routed oracle-optimally. Adding ``\textbf{Accuracy gap on general VQA.}'' as the first item, stated neutrally, will increase reviewer trust in everything else. Note AAAI does not mandate this section; if space is tight, fold it into the conclusion --- but keep the honest accuracy-gap item wherever it lands.}

\textbf{Discrete fidelity set.}
\textsc{VOILA} is designed for the discrete fidelity tiers that characterize real-world tiered storage (captions, thumbnails, compressed JPEGs, full resolution).
Extending the VOI rule to continuous compression schemes is a natural next step that would replace the piecewise policy with a continuous acquisition function.

\textbf{Per-VLM calibration.}
The gradient-boosted regressor is calibrated on correctness labels from a specific VLM.
Collecting these labels is lightweight---3{,}000--12{,}000 examples per configuration suffice---and needs to be done only once offline per VLM--dataset pair.
Cross-model transfer and few-shot calibration are promising directions for model-as-a-service settings.

\textbf{English-language features.}
The current TF-IDF vocabulary and question-type rules target English VQA.
The structural features (question-type flags, lexical patterns) are language-portable, and multilingual extension via language-agnostic encoders is straightforward given labeled data.

\textbf{Synthetic LoCoMo annotations.}
\providecommand{\locomotableref}{\ref{tab:locomo_ultracompact_appendix}}
LoCoMo questions and reference answers are generated by GPT-4o; we use this dataset as an out-of-distribution stress test rather than a primary benchmark, and the strong zero-shot transfer results (Table~\locomotableref) suggest the approach is robust to annotation noise.

\textbf{Cost model grounding.}
The retrieval cost function is derived from empirical file-size ratios and standard tier-latency profiles.
Adapting to a specific deployment requires only re-measuring file sizes and tier latencies---the VOI decision rule is otherwise unchanged.

\textbf{Scope: image VQA.}
\textsc{VOILA} is evaluated on image-based VQA; the VOI framework is modality-agnostic, and applying pre-retrieval fidelity selection to video (temporal sampling), audio (bitrate), and document retrieval (chunk granularity) is a direct extension.
